\section{Elementary maps and elementary substructures}
\label{sec:maps-elementarity-03}

\begin{definition}[Elementary map]
Let \(A\subseteq M\).  A map \(h:A\to N\) is \emph{elementary} iff, for every
\(\Lang\)-formula \(\varphi(x_1,\ldots,x_n)\) and every
\(a_1,\ldots,a_n\in A\),
\[
  \mathcal M\entails\varphi(a_1,\ldots,a_n)
  \quad\Longleftrightarrow\quad
  \mathcal N\entails\varphi(h(a_1),\ldots,h(a_n)).
\]
When \(A=M\), the map is an \emph{elementary embedding}.
\end{definition}

\begin{definition}[Elementary substructure]
If \(\mathcal M\subseteq\mathcal N\) and the inclusion map is elementary,
then \(\mathcal M\) is an \emph{elementary substructure} of \(\mathcal N\),
written \(\mathcal M\preccurlyeq\mathcal N\).  The structure
\(\mathcal N\) is then an \emph{elementary extension} of \(\mathcal M\).
\end{definition}

\begin{remark}
The symbols \(\subseteq\) and \(\preccurlyeq\) express different claims.
The first concerns the interpretations of the signature symbols; the second
concerns all formulas with parameters from the smaller domain.
\end{remark}

\begin{theorem}[Tarski--Vaught test]
\label{thm:tarski-vaught}
Let \(\mathcal M\subseteq\mathcal N\).  Then
\(\mathcal M\preccurlyeq\mathcal N\) iff, for every formula
\(\varphi(y,\bar x)\) and every finite tuple \(\bar a\) from \(M\),
\[
  \mathcal N\entails\exists y\,\varphi(y,\bar a)
  \quad\Longrightarrow\quad
  \mathcal N\entails\varphi(b,\bar a)
  \text{ for some }b\in M.
\]
\textit{Status: complete proof.}
\end{theorem}

This restates \cref{thm:tarski-vaught-fol} in the map-based framework of
this chapter; the proof is repeated so that the present part remains
self-contained.

\begin{proof}
Suppose first that \(\mathcal M\preccurlyeq\mathcal N\) and that the
displayed existential formula is satisfied in \(\mathcal N\).  By
elementarity it is satisfied in
\(\mathcal M\), so a witness \(b\in M\) exists.  Elementarity again gives
\(\mathcal N\entails\varphi(b,\bar a)\).

Conversely, assume the witness condition.  We prove by induction on formulas
\(\psi(\bar x)\) that, for every \(\bar a\) from \(M\),
\[
  \mathcal M\entails\psi(\bar a)
  \quad\Longleftrightarrow\quad
  \mathcal N\entails\psi(\bar a).
\]
Atomic formulas agree because \(\mathcal M\subseteq\mathcal N\), using
\cref{lem:term-preservation}.  Boolean connectives follow from the induction
hypothesis.  For the existential step, if
\(\mathcal M\entails\exists y\,\varphi(y,\bar a)\), a witness in \(M\) and the
induction hypothesis give the same existential statement in \(\mathcal N\).
In the reverse direction, the assumed witness condition supplies \(b\in M\)
with \(\mathcal N\entails\varphi(b,\bar a)\); the induction hypothesis yields
\(\mathcal M\entails\varphi(b,\bar a)\).  For a universal formula
\(\forall y\,\varphi(y,\bar x)\): if
\(\mathcal N\entails\forall y\,\varphi(y,\bar a)\), then for each
\(b\in M\) the instance \(\varphi(b,\bar a)\) is satisfied in
\(\mathcal N\), and the induction hypothesis transfers it to
\(\mathcal M\).  Conversely, if
\(\mathcal N\notentails\forall y\,\varphi(y,\bar a)\), the satisfaction
clauses give \(\mathcal N\entails\exists y\,\neg\varphi(y,\bar a)\); the
witness condition supplies \(b\in M\) with
\(\mathcal N\entails\neg\varphi(b,\bar a)\), the induction hypothesis for
\(\varphi\) gives \(\mathcal M\entails\neg\varphi(b,\bar a)\), and hence
\(\mathcal M\notentails\forall y\,\varphi(y,\bar a)\).
\end{proof}

\begin{example}[Applying the test]
Let \(K\subseteq L\) be algebraically closed fields in the language of rings.
The later quantifier-elimination theorem for algebraically closed fields
implies \(K\preccurlyeq L\).  Without that theorem, closure under the ring
operations establishes only \(K\subseteq L\); algebraic closure alone does
not yield the Tarski--Vaught witness condition without further argument.
\end{example}

